The transverse energy from soft particle production in the underlying event associated with a high transverse energy jet is investigated in p+p collisions at center-of-mass energy of 200 GeV using the sPHENIX detector. Understanding the interplay between jets from the collision’s hard scattering and soft particle emission in the region azimuthally transverse to the event’s hardest jet can inform our understanding of contributions to jet events in proton-proton collisions from multiple parton interactions as well as final and initial state radiation. The sPHENIX calorimeter system, featuring both electromagnetic and hadronic calorimeters with full azimuthal coverage and wide rapidity acceptance, is well suited to make high-resolution measurements of jets and energy measurements of the underlying event transverse to the hard scattering in jet events. This note reports results of the underlying event as a function of jet properties using data collected by sPHENIX from RHIC Run 24 p+p collisions at center-of-mass energy of 200 GeV. 